% Part: turing-machines
% Chapter: undecidability
% Section: halting-problem

\documentclass[../../../include/open-logic-section]{subfiles}

\begin{document}

\olfileid{tur}{und}{hal}
\olsection{مسألة التوقف}

\begin{explain}
افترض أننا ثبتنا تعدادًا ما لأوصاف آلات تورنغ. فتحصل كل آلة تورنغ
بذلك على \emph{فهرس}: موضعها في تعداد أوصاف آلات تورنغ $M_1$، و$M_2$،
و$M_3$،~\dots.

نعلم أنه لا بد من وجود دوال غير قابلة للحساب بآلة تورنغ: فمجموعة
أوصاف آلات تورنغ، ومن ثم مجموعة آلات تورنغ، !!{enumerable}، أما مجموعة
جميع الدوال من $\Nat$ إلى~$\Nat$ فليست كذلك. لكننا نستطيع أيضًا إيجاد
أمثلة معينة لدوال غير قابلة للحساب. ومن هذه الدوال دالة التوقف.
\end{explain}

\begin{defn}[دالة التوقف]
 تُعرّف \emph{دالة التوقف}~$h$ كما يأتي:
\[
h(e,n) =
\begin{cases}
  \text{0} & \text{إذا لم تتوقف الآلة~$M_e$ عند الدخل $n$} \\
  \text{1} & \text{إذا توقفت الآلة~$M_e$ عند الدخل $n$}
\end{cases}
\]
\end{defn}

\begin{defn}[مسألة التوقف]
\emph{مسألة التوقف} هي مسألة تحديد ما إذا كانت آلة تورنغ~$M_e$ تتوقف
عند دخل مؤلف من~$n$ من العلامات (وذلك لأي $e$ و~$n$).
\end{defn}

\begin{explain}
نبين أن~$h$ غير قابلة للحساب بآلة تورنغ ببيان أن دالة مرتبطة بها، هي
$s$، غير قابلة للحساب بآلة تورنغ. ويعتمد هذا البرهان على حقيقة أن كل
ما تستطيع آلة تورنغ حسابه تستطيع آلة تورنغ منضبطة حسابه
(\olref[mac][dis]{sec})، وعلى إمكان وصل آلتي تورنغ لإنتاج آلة واحدة
(\olref[mac][cmb]{sec}).
\end{explain}

\begin{defn} تُعرّف الدالة~$s$ كما يأتي:
\[
s(e) =
\begin{cases}
  \text{0} & \text{إذا لم تتوقف الآلة~$M_e$ عند الدخل $e$} \\
  \text{1} & \text{إذا توقفت الآلة~$M_e$ عند الدخل $e$}
\end{cases}
\]
\end{defn}

\begin{lem}
الدالة~$s$ غير قابلة للحساب بآلة تورنغ.
\end{lem}

\begin{proof}
نفترض، طلبًا للتناقض، أن الدالة~$s$ قابلة للحساب بآلة تورنغ. عندئذ
توجد آلة تورنغ~$S$ تحسب~$s$. ويمكننا أن نفترض، من غير إخلال بالعموم،
أن $S$ حين تتوقف تكون ماسحةً الخانة الأولى (أي إنها منضبطة). ويمكن
«وصل» هذه الآلة بآلة أخرى~$J$ تتوقف إذا بدأت عند الدخل~$0$ (أي إذا
قرأت $\TMblank$ في الحالة الابتدائية وهي تمسح الخانة الواقعة إلى يمين
رمز نهاية الشريط)، وتتيه نحو اليمين فلا تتوقف خلاف ذلك. إن
$S\concat J$، وهي الآلة الناتجة من وصل $S$ بـ~$J$، آلة تورنغ؛ ومن ثم
فهي $M_e$ لبعض~$e$ (أي إنها تظهر في موضع ما من التعداد). شغّل $M_e$
على دخل مؤلف من $e$ من علامات~$\TMstroke$. عندئذ يوجد احتمالان: إما
أن تتوقف $M_e$ وإما ألا تتوقف.
\begin{enumerate}
\item افترض أن $M_e$ تتوقف عند دخل مؤلف من $e$ من علامات~$\TMstroke$.
  عندئذ $s(e)=1$. ومن ثم تتوقف $S$، حين تبدأ عند~$e$، وعلى الشريط
  علامة~$\TMstroke$ واحدة خرجًا. ثم تبدأ $J$ وعلى الشريط علامة
  $\TMstroke$. وفي هذه الحال لا تتوقف $J$. لكن $M_e$ هي الآلة
  $S\concat J$، ولذلك ينبغي أن تفعل بالضبط ما تفعله $S$ متبوعةً بـ~$J$
  (أي أن تتيه في هذه الحال نحو اليمين ولا تتوقف أبدًا). ومن ثم لا يمكن
  أن تتوقف $M_e$ عند دخل مؤلف من $e$ من علامات~$\TMstroke$.

\item افترض الآن أن $M_e$ لا تتوقف عند دخل مؤلف من $e$ من علامات
  $\TMstroke$. عندئذ $s(e)=0$، وتتوقف $S$، حين تبدأ عند الدخل~$e$،
  وشريطها خال. وتتوقف $J$ فورًا حين تبدأ على شريط خال. ومرة أخرى تفعل
  $M_e$ ما تفعله $S$ متبوعةً بـ~$J$، ولذلك لا بد أن تتوقف $M_e$ عند
  دخل مؤلف من $e$ من علامات~$\TMstroke$.
\end{enumerate}
نصل في كلتا الحالتين إلى تناقض مع افتراضنا. وهذا يبين أنه لا يمكن أن
توجد آلة تورنغ~$S$: فالدالة~$s$ غير قابلة للحساب بآلة تورنغ.
\end{proof}

\begin{thm}[استحالة حل مسألة التوقف]
\ollabel{thm:halting-problem} مسألة التوقف غير قابلة للحل، أي إن
الدالة~$h$ غير قابلة للحساب بآلة تورنغ.
\end{thm}

\begin{proof}
افترض أن $h$ قابلة للحساب بآلة تورنغ~$H$. عندئذ يمكننا استعمال $H$
لبناء آلة تورنغ تحسب~$s$: نصنع أولًا نسخة من الدخل (تفصل بينها وبين
الأصل علامة~$\TMblank$)، ثم نعود إلى البداية ونشغّل~$H$. ومن الواضح
أننا نستطيع صنع آلة تنجز المهمة الأولى (انظر
\cref{tur:mac:dis:prob:copier})، ولو وجدت $H$ لاستطعنا «وصلها» بآلة
النسخ هذه لنحصل على آلة جديدة تحدد ما إذا كانت $M_e$ تتوقف عند
الدخل~$e$، أي تحسب~$s$. لكننا بينا بالفعل استحالة وجود آلة كهذه. ومن
ثم تكون $h$ أيضًا غير قابلة للحساب بآلة تورنغ.
\end{proof}

\begin{prob}
مسألة التوقف الثلاثي (3-Halt) هي مسألة إعطاء إجراء قرار يحدد ما إذا
كانت آلة تورنغ مختارة اعتباطيًا تتوقف عند دخل مؤلف من ثلاث علامات
$\TMstroke$ على شريط خال فيما عدا ذلك. أثبت أن مسألة التوقف الثلاثي
غير قابلة للحل.
\end{prob}

\begin{prob}
أثبت أنه إذا كانت مسألة التوقف قابلة للحل لأزواج آلة تورنغ ودخل
$M_e$ و~$n$ التي تحقق $e\neq n$، فهي قابلة للحل أيضًا في الحالات التي
تحقق $e=n$.
\end{prob}

\begin{prob}
برهنا أن مسألة التوقف غير قابلة للحل حين يكون الدخل عددًا~$e$ يعيّن
آلة تورنغ~$M_e$ بواسطة تعداد جميع آلات تورنغ. فماذا لو أتحنا استعمال
أوصاف آلات تورنغ الواردة في \olref[tur][und][enu]{sec} دخلًا مباشرة؟
هل يمكن أن توجد آلة تورنغ تقرر مسألة التوقف لكنها تأخذ أوصاف آلات
تورنغ دخلًا بدلًا من الفهارس؟ وضح لماذا يمكن ذلك أو لا يمكن.
\end{prob}

\begin{prob} أثبت أن الدالة \emph{الجزئية}~$s'$ المعرّفة كما يأتي:
  \[
  s'(e) =
  \begin{cases}
    \text{1} & \text{إذا توقفت الآلة~$M_e$ عند الدخل $e$}\\
    \text{غير معرّفة} & \text{إذا لم تتوقف الآلة~$M_e$ عند الدخل $e$}
  \end{cases}
  \]
  \emph{قابلة} للحساب بآلة تورنغ.
\end{prob}

\end{document}
